\documentclass[a4paper,dvipdfmx]{jsreport}

\usepackage{url}
\usepackage{listings,jlisting}
\usepackage{tcolorbox}
\usepackage{color}
\usepackage{otf}
\usepackage{hyperref}
\usepackage{pxjahyper}

\title{理工系学生のための\LaTeXe 入門ガイドブック}
\author{岩下 雄一郎}

\lstset{
    basicstyle={\ttfamily\small}, %書体の指定
    %language={C},
    commentstyle=\textit,
    keywordstyle=\bfseries,
    frame=single, %フレームの指定
    framesep=10pt, %フレームと中身（コード）の間隔
    breaklines=true, %行が長くなった場合の改行
    numbers=left,
    numbersep={14pt},%行番号と本文との間隔（デフォルトは10pt)
    numberstyle={\sffamily\scriptsize},
    lineskip=-0.5ex, %行間の調整
    tabsize=2, %Tabを何文字幅にするかの指定
    breakindent = 0pt,
    showstringspaces=false
}

\begin{document}

\maketitle
\tableofcontents

\newpage

\part{\LaTeX 環境の構築}
\chapter{\TeX\ Liveのインストール}
\TeX\ Liveは\LaTeX のデストリビューションです．これをインストールすることで\LaTeX 

\section{\LaTeXe 美文書作成入門 を用いたインストール}
\begin{figure}[hbtp]
    \centering
    \includegraphics[scale=0.18]{fig_bibunsho.jpg}
    \caption{［改訂第8版］\LaTeXe 美文書作成入門}
    \label{fig:bibunsho}
\end{figure}

図\ref{fig:bibunsho}は
\begin{quote}
    ［改訂第8版］LaTeX2ε美文書作成入門／奥村晴彦，黒木裕介　著／ISBN 978-4-297-11712-2
\end{quote}
です．こちらはp\LaTeXe を中心として\LaTeX を解説した本で，\LaTeX の主要な機能がほとんど書かれているため，\LaTeX のバイブルのように使うことができます．さらに，付属DVD-ROMを用いて簡単に\TeX\ Live（\LaTeX の実行環境）をインストールできる（WindowsまたはmacOS）ため，購入をおすすめします．

\section{\TeX\ Live公式サイトからインストール}

\TeX\ Liveの公式サイトから無料で\TeX\ Liveをインストールできます．最新版は\TeX\ Live 2021です．インストール方法は下記Webサイトを参照してください．
\begin{description}
    \item[\TeX\ Live 全般] \url{https://texwiki.texjp.org/?TeX%20Live}
    \item[Windows] \url{https://texwiki.texjp.org/?TeX%20Live%2FWindows}
    \item[macOS] \url{https://texwiki.texjp.org/?TeX%20Live%2FMac}
    \item[Linux] \url{https://texwiki.texjp.org/?Linux#texlive}
\end{description}

\chapter{Visual Studio Codeの設定}

本章では\TeX Live がインストール済みである前提でVisual Studio Codeの環境構築方法を解説します．

\section{Visual Studio Code}
Visual Studio Code（VS Code）は無料のコードエディタであり，様々な拡張機能が用意されています．

\section{Visual Studio Codeのインストール}
Visual Studio Codeは
\begin{quote}
    \url{https://code.visualstudio.com}
\end{quote}
からダウンロードできます．このリンクを開き，画面左側のDownloadボタンをクリックするとダウンロードが開始されます．以降はインストーラの指示に従ってコンピュータにインストールしてください．

\section{拡張機能のインストール}
VS Codeを開いてください．

\begin{enumerate}
    \item 左端のバーのExtensionsアイコンを選択．
    \item 検索欄に拡張機能名を入力．
    \item 拡張機能を選択し，Installをクリック．
\end{enumerate}

\LaTeX を書くうえで便利なパッケージを以下に挙げます．

\subsection{Japanese Language Pack for Visual Studio Code}
VS CodeのUIの日本語化パッケージです．Japanese Language Packと調べると出てきます．

\subsection{LaTeX Workshop}
\LaTeX のコンパイル，自動補間などをしてくれるパッケージです．LaTeX Workshopと検索．

\section{LaTeX Workshopの設定}
Visual Studio Codeの拡張機能の設定はsettings.jsonというファイルを通して行えます．

\subsection{settings.jsonを開く}
\begin{enumerate}
    \item ショートカットキー Ctrl\UTF{21E7}P（macOSでは\UTF{21E7}\UTF{2318}P）またはF1でコマンドパレットを表示する．
    \item settingsまたは設定と入力する．
    \item 「基本設定: 設定(JSON)を開く」を選択．
\end{enumerate}

\subsection{settings.jsonの中身}
以下のコードをコピーしてsettings.jsonに貼り付けるか，GitHubにアップロードしたので，そこからダウンロードしてください．

\begin{lstlisting}
{
"latex-workshop.latex.tools": [
    {
        "name": "ptex2pdf",
        "command": "ptex2pdf",
        "args": [
            "-l",
            "-ot",
            "-kanji=utf8",
            "-interaction=nonstopmode",
            "-synctex=1 -file-line-error",
            "%DOCFILE%.tex"
        ]
    }
],
"latex-workshop.latex.recipes": [
    {
        "name": "pLaTeX2e",
        "tools": [
            "ptex2pdf"
        ]
    }
],
"latex-workshop.view.pdf.viewer": "tab",
"editor.renderControlCharacters": true,
"editor.accessibilitySupport": "off"
}
\end{lstlisting}

\part{\LaTeX の基本}

\chapter{\LaTeX とは}
\TeX\footnote{日本では「テフ」や「テック」と呼ばれる．}はスタンフォード大学のDonald E. Knuth教授によって，数学やその他の技術的な文書の組版用に設計されたコンピュータ言語です．\TeX は無料であり，数式を簡単に美しく記述できるため，世界中の多くの研究者によって使われています．

\LaTeX\footnote{呼び方は「ラテフ」や「ラッテック」など．\LaTeX と表記できない場合はLaTeXと書くこともできる．\TeX は本来は別の数式組版ソフトウェアだが，近年は\TeX が\LaTeX を指していることが多い．}は素晴らしい文書作成ツールであり，簡単に綺麗な文書レイアウトを構成できます．しかし，私たちがよく使うMicrosoft WordやApple Pagesとは扱いが全く異なります．普段使っているこれらの文書作成ソフトはWYSIWYG (What You See Is What You Get) ，すなわち，ディスプレイに表示される内容と印刷結果が一致するように表現されたプログラムです．これは直感的に文書を編集できますが，文書構造が曖昧であり，データサイズが大きくなるとメモリを消費してもたつきやすくなります．それに対し，\LaTeX はWYSIWYM (what you see is what you mean) ，すなわち，ユーザは文書内容のみに集中すればよく，印刷物のフォーマットに合わせるのは全てコンパイラが行ってくれます．本章では，\LaTeX による文書作成の基本を解説します．


\chapter{\LaTeX を使った文書作成}

本章では，\LaTeX 初心者の方がすぐに\LaTeX で文書を作成できるようにします．各コマンドの詳細は次章以降で学んでください．それでは，\LaTeX で文書を作成するのがどのような流れで行われるのか体験してみましょう．文字数が多くなってしまいましたが，内容は簡単です．

\section{はじめての\LaTeX による文書作成}
エディタ（指定しません）で新規ファイルを作成し，次のコードを入力してください．
\begin{lstlisting}[language=TeX]
\documentclass{jsarticle}
\begin{document}

Hello, \TeX!
\[ \int dx = x + C. \]

\end{document}
\end{lstlisting}
ここで，次の事柄に注意してください．
\begin{itemize}
    \item 全て直接入力（いわゆる半角文字）で打ち込む
    \item 大文字と小文字を区別（例：\verb|\TeX|と\verb|\Tex|は異なる）
    \item 各行の最後でEnterキーを押して改行．スペースキーではダメ．
    \item 半角の\verb|\|はWindowsでは\verb|¥|のこと．
\end{itemize}

入力し終えたら，ファイル名\footnote{.texファイルのファイル名は英数字が望ましい}を「sample.tex」として保存しましょう．保存場所はホームディレクトリが望ましいです．

それでは，作成した\TeX ファイルをPDF化して行きます．その際，TeX ShopやVS Codeのコンパイル，ビルド機能を使ってボタンをクリックするだけでPDF化するのが一般的ですが，ここでは，\TeX ファイルをPDF化する際の過程を理解するために，コマンドを実行することでPDF化します．

コマンドプロンプト（Windows）かターミナル（macOS，Linux）で次のコマンドを実行してください．これを実行すると，作成した\verb|.tex|ファイルが\textsf{DVIファイル}\footnote{DVIファイル（\textbf{d}e\textbf{v}ice-\textbf{i}ndependent file format，デバイス独立フォーマット）とは，文書のレイアウトなどをデバイスに依存しない形で保存したバイナリファイルです．}（拡張子は\verb|dvi|）に変換されます．この作業は\textgt{コンパイル}と呼ばれます．

\begin{lstlisting}[language=sh]
platex sample.tex
\end{lstlisting}
実行時に以下のメッセージが画面に表示されると思います．ここには，\verb|sample.dvi|が生成されたと書かれています．また，\verb|.tex|ファイルがp\LaTeXe の文法に則っていない場合はメッセージにエラーが表示されます．
\begin{lstlisting}[language=sh]
This is e-pTeX, Version 3.14159265-p3.8.3-191112-2.6 (utf8.euc) (TeX Live 2020) (preloaded format=platex)
restricted \write18 enabled.
entering extended mode
(./sample.tex
pLaTeX2e <2021-06-01>+2 (based on LaTeX2e <2021-06-01> patch level 1)
L3 programming layer <2021-08-27>
(/Applications/TeXLive/Library/texlive/2021/texmf-dist/tex/platex/jsclasses/jsa
rticle.cls
Document Class: jsarticle 2021/06/28 jsclasses (okumura, texjporg)

(/Applications/TeXLive/Library/texlive/2021/texmf-dist/tex/platex/jsclasses/jsl
ogo.sty))
(/Applications/TeXLive/Library/texlive/2021/texmf-dist/tex/latex/l3backend/l3ba
ckend-dvipdfmx.def (|extractbb --version))
No file sample.aux.
[1] (./sample.aux) )
Output written on sample.dvi (1 page, 380 bytes).
Transcript written on sample.log.
\end{lstlisting}

これで，DVIファイルが生成されました．しかし，これはバイナリファイルなので，私たちが読むことができる形式で出力しなければなりません．ここでは，DVIファイルをPDFファイルに変換します．\footnote{PDFを出力するのが一般的ですが，Postscriptファイルに変換してプレビュー・印刷することも可能です．}

次に，タイプセット（コンパイル）ボタンを押してください．
\begin{itemize}
    \item TeXShopなら左上の\fbox{タイプセット}ボタン．
    \item TeXworksなら$\triangleright$ボタン
    \item VSCodeなら緑色の$\triangleright$ボタン
\end{itemize}
を押してください．

次の内容がPDFに出力されるはずです．

\begin{tcolorbox}
    Hello, \TeX!
    \[ \int dx = x + C. \]
\end{tcolorbox}

ここでは，次のような処理が行われています．
\begin{enumerate}
    \item articleというドキュメントクラス（文書の種類）で文書を作成する．（1行目）
    \item 文書を作成せよ．（2行目）
    \item 文書の一段落目に
        \begin{tcolorbox}
            Hello, \TeX!
            \[ \int dx = x + C. \]
        \end{tcolorbox}
        を出力せよ．（4，5行目）
    \item 文書を終了せよ．（7行目）
\end{enumerate}

このように，\LaTeX では全て自分でレイアウト・文書構造をコマンドを用いて設定します．これがワープロソフトと大きく異なる点ではないでしょうか．

次ページから\LaTeX で文書を作成するための道具を紹介して行きます．

\chapter{文書作成の基礎}
\section{\LaTeX 文書の型}
\subsection{文書クラス}
Wordでは新規ファイルを開いていきなり本文が書けますが，\LaTeX では違います．まず，\textgt{文書（ドキュメント）クラス}を設定する必要があります．これは，どのテンプレートや用紙サイズを用いるかなど，文書のレイアウトの大枠を決めるものです．次のようなコマンドを書きます．
\begin{quote}
    \verb|\documentclass[オプション]{文書クラス名}|
\end{quote}

文書クラスには主に表\ref{tab:documentclass}に示すものがあります．基本的には和文（新・横）と書かれているものを使用してください．jsclassesの最新バージョンで使用可能です．uをつけたものはup\LaTeX 用です．作成する文書の形式に合うものを選び，文書クラス名の欄に入力してください．
\begin{table}[h]
    \centering
    \caption{文書クラス}
    \label{tab:documentclass}
    \begin{tabular}{|c|llll|}
        \hline
        用途 & 欧文 & 和文（旧・横) & 和文（旧・縦） & 和文（新・横）\\ \hline
        論文・レポート & \verb|article| & \verb|(u)jarticle| & \verb|(u)tarticle| & \verb|jsarticle| \\
        長い報告書 & \verb|report| & \verb|(u)jreport| & \verb|(u)treport| & \verb|jsreport| \\
        本 & \verb|book| & \verb|(u)jbook| & \verb|(u)tbook| & \verb|jsbook| \\ \hline
    \end{tabular}
\end{table}

次に，オプション欄では用紙サイズや本文の文字サイズを設定できます．なお，オプションを書かないこと（\verb|\documentclass{文書クラス名}|と入力）もでき，この場合はデフォルト設定となります．デフォルトでは用紙サイズA4（縦），本文の文字サイズ10ptです．

\subsection{文書の開始と終了}
文書クラスを設定しただけではまだ文章を書こうとしてもタイプセットできません．どこからどこまでが文書か設定する必要があります．次のように，本文を\verb|\begin{document}|と\verb|\end{document}|で挟みます．

\begin{lstlisting}[language=TeX]
\documentclass{jsarticle}
\begin{document}
本文をここに書きます．
\end{document}
\end{lstlisting}

これをタイプセットすると次のようになります．
\begin{tcolorbox}
本文をここに書きます．
\end{tcolorbox}

\subsection{注意}
以上で本文を出力できるようになりましたが，どんな文字でも出力できるわけではありません．後ほど説明しますので，出力できない文字があっても気にしないでください．

\section{文書のレイアウト}
\LaTeX ではWordと異なり，タイトルや章構成を明示的に記述する必要があります．それらは\textgt{コマンド}により設定します．
\subsection{表題}
\LaTeX では表題をコマンドにより表示することができます．デフォルトでは表題の内容は\textgt{タイトル}，\textgt{著者}，\textgt{日付}です．\verb|\documentclass{...}|と\verb|\begin{document}|の間の部分を\textgt{プリアンブル}と呼びます．ここに次のコマンドを書いてください．
\begin{quote}
    \verb|\title{文書タイトル}|\\
    \verb|\autor{著者}|\\
    \verb|\date{日付}|
\end{quote}

そして，タイトルを表示したい部分に
\begin{quote}
    \verb|\maketitle|
\end{quote}
と入力してください．\verb|\maketitle|が無ければ表題が表示されません．

表題を表示する場合，\verb|\title|と\verb|\author|コマンドは必須です．\verb|\date|コマンドは任意ですが，書かない場合はタイプセットした日の日付が出力されます．また，これらのコマンドの引数の書き方に指定はありません．著者や日付を出力したくない場合は\verb|\author{}|や\verb|\date{}|などと括弧の中身を空欄にしてください．

\subsection{目次}
任意で目次を表示することができます．目次を表示したい箇所に
\begin{quote}
    \verb|\tableofcontents|
\end{quote}
と入力してください．

\subsection{見出し}
\LaTeX ではWordと異なり，見出しを明示的に書きます．文書クラスによって使えない見出しがあるので表で確認してください．

\begin{table}[ht]
    \centering
    \caption{見出し}
    \label{tb:title}
    \begin{tabular}{|c|l|ccc|}
        \hline
        説明 & コマンド & \verb|jsbook| & \verb|jsreport| & \verb|jsarticle| \\ \hline
        部 & \verb|\part{部}| & ◯ & ◯ & × \\
        章 & \verb|\chapter{章}| & ◯ & ◯ & ◯ \\
        節 & \verb|\section{節}| & ◯ & ◯ & ◯ \\
        小節 & \verb|\subsection{小節}| & ◯ & ◯ & ◯ \\
        小々節 & \verb|\subsubsection{小々節}| & ◯ & ◯ & ◯ \\ \hline
    \end{tabular}
\end{table}

\section{注釈／コメント}
\LaTeX では注釈の先頭に\verb|%|を書くことで注釈を書くことができます．例えば，次のように入力します．
\begin{lstlisting}
\documentclass{jsarticle}
% コメントをここに書きます．
\begin{document}
本文をここに書きます．% 本文
\[ \int dx = x + C. \] % 数式
\end{document}
\end{lstlisting}
これをタイプセットすると次のようになります．
\begin{tcolorbox}
本文をここに書きます．% 本文
\[ \int dx = x + C. \] % 数式
\end{tcolorbox}

\section{文字の扱い}
\subsection{書体}
ここでは，パッケージなしのデフォルトの状態で使える書体を紹介します．
\subsubsection{和文}
和文では\,\fbox{明朝体}\,と\,\fbox{\textgt{ゴシック体}}\,の2種類の書体が使えます．明朝体はデフォルトの和文書体です\footnote{明朝体を指定するコマンドも存在しますがここでは省略します．}．一方，ゴシック体を使用する際は\verb|\textgt{...}|コマンドで指定する必要があります．
\subsubsection{欧文}
欧文では
\section{注意事項}

\begin{thebibliography}{00}
    \bibitem{texlive-ja} Karl Berry 編, 朝倉卓人 (wtsnjp) 訳, ``\TeX Live ガイド 2021,'' 2021年3月.
\end{thebibliography}
\end{document}